\documentclass{zhvt-classic}

% 使用本地源流明体（GenRyuMin2 TW）按文件路径加载，不依赖系统字体名
% 版框（内框／界栏／书口侧栏／鱼尾）改用与版心印章同一印泥棕色：
%   印章图片 yshs.png 实色 = RGB(135,68,52) = HTML 874434
\usepackage{xcolor}
\usepackage{tikz}
% sealbrown 双注册：xcolor 供 TikZ/pgf，L3 供 cls 版框
\definecolor{sealbrown}{HTML}{874434}
\ExplSyntaxOn
\color_set:nnn { sealbrown } { HTML } { 874434 }
\ExplSyntaxOff

\zhvtset{
  grid color       = sealbrown,
  main font        = GenRyuMin2TW-SB,
  main font path   = C:/Users/Administrator/texmf/fonts/opentype/public/GenRyuMin2TW/,
jiazhu font      = GenRyuMin2TW-M,
  jiazhu font path = C:/Users/Administrator/texmf/fonts/opentype/public/GenRyuMin2TW/,
  % 页码框容纳 4 字：正常页「九十七」+前置空 1 格=4 行恰好；
  % 100+ 页「一百零一」共 5 行超出，末字「一/二/三」被框裁掉（2026-10-04 实测），故放宽到 5
  shukou page number chars = 5,
  % 书口书名框高从默认 6 字减到 4 字，把省出的 2 字高拨给 chapter 框，
  % 使 6 字书名+卷名放进 9 字高的 chapter 框后不再压线（2026-10-05）
  shukou title chars = 4,
}

\title{史记}[廉颇蔺相如列传]
\author{汉·司马迁}
\maker{闲叟漫记·重製}


% ===== 索隐述赞样式（2026-11-08 四次定稿 · 愚大人亲手改定参数）=====
% ⚠ 竖排前提：自定义字体家族必须带 RawFeature={vertical:+vert}；
%   simhei.ttf 加 +vhal 会整体加载失败静默回退，只写 +vert；simfang 同理。
\newCJKfontfamily{\shuzanfont}[Path=C:/Windows/Fonts/,RawFeature={vertical:+vert}]{simhei.ttf}
\newCJKfontfamily{\shuzanbody}[Path=C:/Windows/Fonts/,RawFeature={vertical:+vert}]{simfang.ttf}
% 题头：152x30pt 圆角 6pt 棕底白字；baseline (0,2) 使整框落入界行格子不压版框线；
% 白字节点 x 不补偿、y +8.2pt（像素级墨迹实测 dx/dy≈0，L/R 留白约 20pt）。
\newcommand{\shuzanhead}[1]{%
  \begin{tikzpicture}[baseline={(0,2)}]
    \fill[sealbrown, rounded corners=6pt]
      (-76pt,-15pt) rectangle (76pt,15pt);
    \node[text=white, inner sep=0pt, anchor=center] at (0pt,8.2pt)
      {\shuzanfont\fontsize{23}{23}\selectfont #1};
  \end{tikzpicture}%
}
% ★ 修正书口卷名位置：chapter 框起始位置改回 cls 默认值，避免全局影响。
%   只在每个卷首页通过 \juan 命令里的空命令微调该页位置（安全，不会跨页污染）。
\ExplSyntaxOn
\AtBeginDocument{
  \dim_set:Nn \l__zhvt_shukou_chapter_title_start_dim
    { \l__zhvt_shukou_title_height_dim + \l__zhvt_font_size_dim * 2 }
  \dim_set:Nn \l__zhvt_shukou_page_number_start_dim
    { \l__zhvt_shukou_chapter_title_start_dim + \l__zhvt_shukou_chapter_title_height_dim }
}
\ExplSyntaxOff

\begin{document}

% ★★ 版心（书口）排布覆写：把书名放到鱼尾之下（2026-10-04）
%
%   cls 默认版心自上而下：书名（title 静态框，在鱼尾之上）→ 鱼尾 → 篇名（chapter 动态框）→ 页码
%   本模板改为：          鱼尾（位置不动）→ 书名 → 篇名 → 页码分割线 → 页码
%
%   手法：把 cls 放在 title 静态框里的书名，改放进 chapter 动态框（鱼尾之下）；
%         title 静态框留空。书名取 \l__zhvt_title_str，即 \title{...} 的第一参数，
%         所以换书时只改 \title，不必在此重复写书名。
%
%   ⚠ 三条必须遵守的写法（2026-10-04 实测）：
%     ① 汉字必须用 \hbox{...} 包住。flowfram 会把框内容装进一个宽 14.4mm（= 一个 line_sep
%        = 书口栏宽）的纵向 minipage 并按该宽度折行，而整页靠 PDF /Rotate 90 显示、
%        框内文字本身不变向 —— 不包 \hbox 时「围炉夜话」4 字会被逐字折成 4 行，
%        在版面上就成了横排 4 字并横向溢出版框。
%     ② 竖排下**一个 \hbox 就是一整列**。书名与篇名要同处一列，必须放进同一个 \hbox；
%        分成两个 \hbox 时，后一个会被当成新的一列、甩到隔壁栏去
%        （实测「序」跑到书口栏外侧 14.4mm 的正文列上）。中间的 \quad 在框内是横向空白，
%        显示出来正好是书名与篇名之间的纵向间隔。
%     ③ ★ 内容盒要写成 **与框等宽的定宽盒** `\hbox to \l__zhvt_shukou_chapter_title_height_dim{…\hfil}`。
%        实测：chapter 框是把内容盒沿纵向**居中**摆放的，内容宽度一变，书名就跟着漂：
%        自然宽度的 \hbox（书名 38.4mm）落在 x=157.2mm，加篇名后（57.6mm）落在 x=152.4mm，
%        目录页（无篇名）与正文页的书名就差半格。做成等宽盒后，无论有没有篇名，
%        内容固定从框顶起排，实测两页均落在 x=142.8mm，不再漂。
\ExplSyntaxOn
% 把书口 chapter 框（书名+篇名）起始位置下移 1 字高，避免 4 字卷名顶到鱼尾
% ⚠ cls 的 \__zhvt_calculate_dimensions: 在 \AtBeginDocument 里执行，会重算该 dim；
%   必须把覆写挂到 begin document 之后，否则会被二次覆盖（2026-10-05 实测）。
% (旧方案已移到导言区，直接修改 chapter 框起始位置)
% 卷标题命令 #1=书口篇名(鱼尾) #2=正文卷名（含卷序，卷名内空格用 ~）
%
% ★★ 目录页两项修订（2026-10-05，愚大人指出）★★
%   ① 页码起点统一对齐：
%      原状 —— 目录项写成 \par\hspace{4em}卷名\hspace{4em}页码，页码起点 = 卷名实宽 + 4em。
%      卷名实宽随字数变：「卷一~梁惠王上」6 字 +~、「卷七~离娄上」5 字 +~，
%      差 1 个字位 → 卷七~卷十 4 条的页码整体比其余 10 条高 27.21pt（一个字），看着上下不齐。
%      改法 —— 卷名补全角空格（宽 = 1 字 = 27.215pt）补足到 **7 字符位**再写进目录，
%      14 条页码起点遂全部落在同一条竖线上（起点对齐；末字随页码字数自然参差）。
%      ⚠ 只对目录里的副本 \l_tmpc_tl 补，正文卷名仍用 \l_tmpa_tl 原样 —— 正文像素不受影响。
%   ② 卷名整体上移半字位：前导 \hspace 由 4em 改 3.5em。
%      实测原状目录块在版框内偏下（天头留白 139.6pt / 地脚留白 110.7pt），
%      上移 0.5em = 13.6pt 后上下基本均分，视觉居中。
\NewDocumentCommand { \juan } { mm } {
	\clearpage
	\tl_set:Nn \l__zhvt_chapter_title_tl { #1 }
	\tl_set:Nn \l_tmpa_tl { #2 }
	\tl_set_eq:NN \l_tmpc_tl \l_tmpa_tl
	\int_step_inline:nn { 7 - \str_count:n { #2 } }
	  { \tl_put_right:Nn \l_tmpc_tl { 　 } }
	\tl_set:Nx \l_tmpb_tl { \zhnum { page } }
	\tl_put_left:Nn \l_tmpb_tl { \hspace { 4em } }
	\tl_put_left:Nn \l_tmpb_tl { \l_tmpc_tl }
	\tl_put_left:Nn \l_tmpb_tl { \par \hspace { 3.5em } }
	\addtocontents { toc } { \l_tmpb_tl }
	\skip_horizontal:n { \l__zhvt_font_size_dim * 4 }
	\l_tmpa_tl
	\par
	\thispagestyle { empty }
}
\cs_set_protected:Npn \__zhvt_set_shukou_content: {
	% ★ 鱼尾上方留白（title 静态框）嵌入闲叟钤印：
	%   整页 /Rotate 90，故此处把图用 \rotatebox{90} 转正为竖章；
	%   尺寸卡在版心条幅内（书内约 7mm 宽 × 33.8mm 高，留余量不压鱼尾）。
	\setstaticcontents * { title } {
		\vbox_to_ht:nn { \l__zhvt_baseline_skip_dim } {
			\vskip 14.2pt
			\hbox_to_wd:nn { \l__zhvt_shukou_title_height_dim } {
				\hfil
				\rotatebox [ origin = c ] { 90 } { \includegraphics [ width = 7mm , height = 33.8mm ] { yshs.png } }
				\hfil
			}
		}
	}
% ★ 版心（书口）统一改法（2026-11-16 愚大人定）：18pt＋框内居中＋书口文字同版框色
%   ① 原 27.2pt 时 7 字长篇名（如「廉颇蔺相如列传」）末字溢出 chapter 框、压住页码；
%      改 \fontsize{18pt}{18pt}，前后 \hfil 令字串在框内居中（不随篇名字数漂移）。
%   ② \vbox{\vskip 8.2pt} 定量下移，抵消缩字号后栏宽方向 8.2pt 上偏（像素实测 y 中心 543.3 = 原版）。
%      ⚠ \raisebox / \vspace* 在 flowfram 框内无效；换字号须同步换补偿：16pt→10.1、19pt→7.4。
%   ③ chapter 与 page~number 均置 \color{sealbrown}（HTML 874434），与版框／界行同色；正文仍黑。
%   ⚠ 页码字位、钤印、鱼尾、界行、正文一律未动；五部合订本同此口径。
	\setdynamiccontents * { chapter }
	{ \color{sealbrown} \vbox { \vskip 8.2pt \hbox to \l__zhvt_shukou_chapter_title_height_dim
		{ \hfil \fontsize{18pt}{18pt}\selectfont \l__zhvt_title_str \quad \l__zhvt_chapter_title_tl \hfil } } }
	\setdynamiccontents * { page~number } { \color{sealbrown} \quad \zhnum { page } }
}
\__zhvt_set_shukou_content:
\ExplSyntaxOff

\maketitle

% 单篇，无多卷目录，不渲染目录页

\juan{廉颇蔺相如列传}{卷八十一~廉颇蔺相如列传第二十一}
\hspace*{2em}廉颇者。赵之良将也。赵惠文王十六年。廉颇为赵将伐齐。大破之。取阳晋。拜为上卿。以勇气闻於诸侯。蔺相如者。赵人也。为赵宦者令缪贤舍人。\\
\hspace*{2em}赵惠文王时。得楚和氏璧。秦昭王闻之。使人遗赵王书。原以十五城请易璧。赵王与大将军廉颇诸大臣谋。欲予秦。秦城恐不可得。徒见欺。欲勿予。即患秦兵之来。计未定。求人可使报秦者。未得。宦者令缪贤曰。﹁臣舍人蔺相如可使。﹂王问。﹁何以知之。﹂对曰。﹁臣尝有罪。窃计欲亡走燕。臣舍人相如止臣。曰。﹃君何以知燕王。﹄臣语曰。﹃臣尝从大王与燕王会境上。燕王私握臣手。曰﹁原结友﹂。以此知之。故欲往。﹄相如谓臣曰。﹃夫赵强而燕弱。而君幸於赵王。故燕王欲结於君。今君乃亡赵走燕。燕畏赵。其势必不敢留君。而束君归赵矣。君不如肉袒伏斧质请罪。则幸得脱矣。﹄臣从其计。大王亦幸赦臣。臣窃以为其人勇士。有智谋。宜可使。﹂於是王召见。问蔺相如曰。﹁秦王以十五城请易寡人之璧。可予不。﹂相如曰。﹁秦强而赵弱。不可不许。﹂王曰。﹁取吾璧。不予我城。奈何。﹂相如曰。﹁秦以城求璧而赵不许。曲在赵。赵予璧而秦不予赵城。曲在秦。均之二策。宁许以负秦曲。﹂王曰。﹁谁可使者。﹂相如曰。﹁王必无人。臣原奉璧往使。城入赵而璧留秦。城不入。臣请完璧归赵。﹂赵王於是遂遣相如奉璧西入秦。\\
\hspace*{2em}秦王坐章台见相如。相如奉璧奏秦王。秦王大喜。传以示美人及左右。左右皆呼万岁。相如视秦王无意偿赵城。乃前曰。﹁璧有瑕。请指示王。﹂王授璧。相如因持璧卻立。倚柱。怒发上冲冠。谓秦王曰。﹁大王欲得璧。使人发书至赵王。赵王悉召群臣议。皆曰﹃秦贪。负其强。以空言求璧。偿城恐不可得﹄。议不欲予秦璧。臣以为布衣之交尚不相欺。况大国乎。且以一璧之故逆强秦之驩。不可。於是赵王乃斋戒五日。使臣奉璧。拜送书於庭。何者。严大国之威以修敬也。今臣至。大王见臣列观。礼节甚倨。得璧。传之美人。以戏弄臣。臣观大王无意偿赵王城邑。故臣复取璧。大王必欲急臣。臣头今与璧俱碎於柱矣。﹂相如持其璧睨柱。欲以击柱。秦王恐其破璧。乃辞谢固请。召有司案图。指从此以往十五都予赵。相如度秦王特以诈详为予赵城。实不可得。乃谓秦王曰。﹁和氏璧。天下所共传宝也。赵王恐。不敢不献。赵王送璧时。斋戒五日。今大王亦宜斋戒五日。设九宾於廷。臣乃敢上璧。﹂秦王度之。终不可强夺。遂许斋五日。舍相如广成传。相如度秦王虽斋。决负约不偿城。乃使其从者衣褐。怀其璧。从径道亡。归璧于赵。\\
\hspace*{2em}秦王斋五日後。乃设九宾礼於廷。引赵使者蔺相如。相如至。谓秦王曰。﹁秦自缪公以来二十馀君。未尝有坚明约束者也。臣诚恐见欺於王而负赵。故令人持璧归。间至赵矣。且秦强而赵弱。大王遣一介之使至赵。赵立奉璧来。今以秦之强而先割十五都予赵。赵岂敢留璧而得罪於大王乎。臣知欺大王之罪当诛。臣请就汤镬。唯大王与群臣孰计议之。﹂秦王与群臣相视而嘻。左右或欲引相如去。秦王因曰。﹁今杀相如。终不能得璧也。而绝秦赵之驩。不如因而厚遇之。使归赵。赵王岂以一璧之故欺秦邪。﹂卒廷见相如。毕礼而归之。\\
\hspace*{2em}相如既归。赵王以为贤大夫使不辱於诸侯。拜相如为上大夫。秦亦不以城予赵。赵亦终不予秦璧。\\
\hspace*{2em}其後秦伐赵。拔石城。明年。复攻赵。杀二万人。\\
\hspace*{2em}秦王使使者告赵王。欲与王为好会於西河外渑池。赵王畏秦。欲毋行。廉颇。蔺相如计曰。﹁王不行。示赵弱且怯也。﹂赵王遂行。相如从。廉颇送至境。与王诀曰。﹁王行。度道里会遇之礼毕。还。不过三十日。三十日不还。则请立太子为王。以绝秦望。﹂王许之。遂与秦王会渑池。秦王饮酒酣。曰。﹁寡人窃闻赵王好音。请奏瑟。﹂赵王鼓瑟。秦御史前书曰﹁某年月日。秦王与赵王会饮。令赵王鼓瑟﹂。蔺相如前曰。﹁赵王窃闻秦王善为秦声。请奏盆缻秦王。以相娱乐。﹂秦王怒。不许。于是相如前进缻。因跪请秦王。秦王不肯击缻。相如曰。﹁五步之内。相如请得以颈血溅大王矣。﹂左右欲刃相如。相如张目叱之。左右皆靡。於是秦王不怿。为一击缶。相如顾召赵御史书曰﹁某年月日。秦王为赵王击鲊﹂。秦之群臣曰。﹁请以赵十五城为秦王寿﹂。蔺相如亦曰。﹁请以秦之咸阳为赵王寿。﹂秦王竟酒。终不能加胜於赵。赵亦盛设兵以待秦。秦不敢动。\\
\hspace*{2em}既罢归国。以相如功大。拜为上卿。位在廉颇之右。廉颇曰。﹁我为赵将。有攻城野战之大功。而蔺相如徒以口舌为劳。而位居我上。且相如素贱人。吾羞。不忍为之下。﹂宣言曰。﹁我见相如。必辱之。﹂相如闻。不肯与会。相如每朝时。常称病。不欲与廉颇争列。已而相如出。望见廉颇。相如引车避匿。於是舍人相与谏曰。﹁臣所以去亲戚而事君者。徒慕君之高义也。今君与廉颇同列。廉君宣恶言而君畏匿之。恐惧殊甚。且庸人尚羞之。况於将相乎。臣等不肖。请辞去。﹂蔺相如固止之。曰。﹁公之视廉将军孰与秦王。﹂曰。﹁不若也。﹂相如曰。﹁夫以秦王之威。而相如廷叱之。辱其群臣。相如虽驽。独畏廉将军哉。顾吾念之。强秦之所以不敢加兵於赵者。徒以吾两人在也。今两虎共斗。其势不俱生。吾所以为此者。以先国家之急而後私雠也。﹂廉颇闻之。肉袒负荆。因宾客至蔺相如门谢罪。曰。﹁鄙贱之人。不知将军宽之至此也。﹂卒相与驩。为刎颈之交。\\
\hspace*{2em}是岁。廉颇东攻齐。破其一军。居二年。廉颇复伐齐几。拔之。後三年。廉颇攻魏之防陵。安阳。拔之。後四年。蔺相如将而攻齐。至平邑而罢。其明年。赵奢破秦军阏与下。\\
\hspace*{2em}赵奢者。赵之田部吏也。收租税而平原君家不肯出租。奢以法治之。杀平原君用事者九人。平原君怒。将杀奢。奢因说曰。﹁君於赵为贵公子。今纵君家而不奉公则法削。法削则国弱。国弱则诸侯加兵。诸侯加兵是无赵也。君安得有此富乎。以君之贵。奉公如法则上下平。上下平则国强。国强则赵固。而君为贵戚。岂轻於天下邪。﹂平原君以为贤。言之於王。王用之治国赋。国赋大平。民富而府库实。\\
\hspace*{2em}秦伐韩。军於阏与。王召廉颇而问曰。﹁可救不。﹂对曰。﹁道远险狭。难救。﹂又召乐乘而问焉。乐乘对如廉颇言。又召问赵奢。奢对曰。﹁其道远险狭。譬之犹两鼠斗於穴中。将勇者胜。﹂王乃令赵奢将。救之。\\
\hspace*{2em}兵去邯郸三十里。而令军中曰。﹁有以军事谏者死。﹂秦军军武安西。秦军鼓譟勒兵。武安屋瓦尽振。军中候有一人言急救武安。赵奢立斩之。坚壁。留二十八日不行。复益增垒。秦间来入。赵奢善食而遣之。间以报秦将。秦将大喜曰。﹁夫去国三十里而军不行。乃增垒。阏与非赵地也。﹂赵奢既已遣秦间。卷甲而趋之。二日一夜至。今善射者去阏与五十里而军。军垒成。秦人闻之。悉甲而至。军士许历请以军事谏。赵奢曰。﹁内之。﹂许历曰。﹁秦人不意赵师至此。其来气盛。将军必厚集其阵以待之。不然。必败。﹂赵奢曰。﹁请受令。﹂许历曰。﹁请就鈇质之诛。﹂赵奢曰。﹁胥後令邯郸。﹂许历复请谏。曰。﹁先据北山上者胜。後至者败。﹂赵奢许诺。即发万人趋之。秦兵後至。争山不得上。赵奢纵兵击之。大破秦军。秦军解而走。遂解阏与之围而归。\\
\hspace*{2em}赵惠文王赐奢号为马服君。以许历为国尉。赵奢於是与廉颇。蔺相如同位。\\
\hspace*{2em}後四年。赵惠文王卒。子孝成王立。七年。秦与赵兵相距长平。时赵奢已死。而蔺相如病笃。赵使廉颇将攻秦。秦数败赵军。赵军固壁不战。秦数挑战。廉颇不肯。赵王信秦之间。秦之间言曰。﹁秦之所恶。独畏马服君赵奢之子赵括为将耳。﹂赵王因以括为将。代廉颇。蔺相如曰。﹁王以名使括。若胶柱而鼓瑟耳。括徒能读其父书传。不知合变也。﹂赵王不听。遂将之。\\
\hspace*{2em}赵括自少时学兵法。言兵事。以天下莫能当。尝与其父奢言兵事。奢不能难。然不谓善。括母问奢其故。奢曰。﹁兵。死地也。而括易言之。使赵不将括即已。若必将之。破赵军者必括也。﹂及括将行。其母上书言於王曰。﹁括不可使将。﹂王曰。﹁何以。﹂对曰。﹁始妾事其父。时为将。身所奉饭饮而进食者以十数。所友者以百数。大王及宗室所赏赐者尽以予军吏士大夫。受命之日。不问家事。今括一旦为将。东向而朝。军吏无敢仰视之者。王所赐金帛。归藏於家。而日视便利田宅可买者买之。王以为何如其父。父子异心。原王勿遣。﹂王曰。﹁母置之。吾已决矣。﹂括母因曰。﹁王终遣之。即有如不称。妾得无随坐乎。﹂王许诺。\\
\hspace*{2em}赵括既代廉颇。悉更约束。易置军吏。秦将白起闻之。纵奇兵。详败走。而绝其粮道。分断其军为二。士卒离心。四十馀日。军饿。赵括出锐卒自博战。秦军射杀赵括。括军败。数十万之众遂降秦。秦悉阬之。赵前後所亡凡四十五万。明年。秦兵遂围邯郸。岁馀。几不得脱。赖楚。魏诸侯来救。乃得解邯郸之围。赵王亦以括母先言。竟不诛也。\\
\hspace*{2em}自邯郸围解五年。而燕用栗腹之谋。曰﹁赵壮者尽於长平。其孤未壮﹂。举兵击赵。赵使廉颇将。击。大破燕军於鄗。杀栗腹。遂围燕。燕割五城请和。乃听之。赵以尉文封廉颇为信平君。为假相国。\\
\hspace*{2em}廉颇之免长平归也。失势之时。故客尽去。及复用为将。客又复至。廉颇曰。﹁客退矣。﹂客曰。﹁吁。君何见之晚也。夫天下以市道交。君有势。我则从君。君无势则去。此固其理也。有何怨乎。﹂居六年。赵使廉颇伐魏之繁阳。拔之。\\
\hspace*{2em}赵孝成王卒。子悼襄王立。使乐乘代廉颇。廉颇怒。攻乐乘。乐乘走。廉颇遂奔魏之大梁。其明年。赵乃以李牧为将而攻燕。拔武遂。方城。\\
\hspace*{2em}廉颇居梁久之。魏不能信用。赵以数困於秦兵。赵王思复得廉颇。廉颇亦思复用於赵。赵王使使者视廉颇尚可用否。廉颇之仇郭开多与使者金。令毁之。赵使者既见廉颇。廉颇为之一饭斗米。肉十斤。被甲上马。以示尚可用。赵使还报王曰。﹁廉将军虽老。尚善饭。然与臣坐。顷之三遗矢矣。﹂赵王以为老。遂不召。\\
\hspace*{2em}楚闻廉颇在魏。阴使人迎之。廉颇一为楚将。无功。曰。﹁我思用赵人。﹂廉颇卒死于寿春。\\
\hspace*{2em}李牧者。赵之北边良将也。常居代雁门。备匈奴。以便宜置吏。市租皆输入莫府。为士卒费。日击数牛飨士。习射骑。谨烽火。多间谍。厚遇战士。为约曰。﹁匈奴即入盗。急入收保。有敢捕虏者斩。﹂匈奴每入。烽火谨。辄入收保。不敢战。如是数岁。亦不亡失。然匈奴以李牧为怯。虽赵边兵亦以为吾将怯。赵王让李牧。李牧如故。赵王怒。召之。使他人代将。\\
\hspace*{2em}岁馀。匈奴每来。出战。出战。数不利。失亡多。边不得田畜。复请李牧。牧杜门不出。固称疾。赵王乃复强起使将兵。牧曰。﹁王必用臣。臣如前。乃敢奉令。﹂王许之。\\
\hspace*{2em}李牧至。如故约。匈奴数岁无所得。终以为怯。边士日得赏赐而不用。皆原一战。於是乃具选车得千三百乘。选骑得万三千匹。百金之士五万人。彀者十万人。悉勒习战。大纵畜牧。人民满野。匈奴小入。详北不胜。以数千人委之。单于闻之。大率众来入。李牧多为奇陈。张左右翼击之。大破杀匈奴十馀万骑。灭襜褴。破东胡。降林胡。单于奔走。其後十馀岁。匈奴不敢近赵边城。\\
\hspace*{2em}赵悼襄王元年。廉颇既亡入魏。赵使李牧攻燕。拔武遂。方城。居二年。庞暖破燕军。杀剧辛。後七年。秦破杀赵将扈辄於武遂。斩首十万。赵乃以李牧为大将军。击秦军於宜安。大破秦军。走秦将桓齮。封李牧为武安君。居三年。秦攻番吾。李牧击破秦军。南距韩。魏。\\
\hspace*{2em}赵王迁七年。秦使王翦攻赵。赵使李牧。司马尚御之。秦多与赵王宠臣郭开金。为反间。言李牧。司马尚欲反。赵王乃使赵葱及齐将颜聚代李牧。李牧不受命。赵使人微捕得李牧。斩之。废司马尚。後三月。王翦因急击赵。大破杀赵葱。虏赵王迁及其将颜聚。遂灭赵。\\
\hspace*{2em}太史公曰。知死必勇。非死者难也。处死者难。方蔺相如引璧睨柱。及叱秦王左右。势不过诛。然士或怯懦而不敢发。相如一奋其气。威信敌国。退而让颇。名重太山。其处智勇。可谓兼之矣。\\

% ===== 索隐述赞（司马贞《史记索隐》）· 统一样式 =====
% 题头「索隐述赞曰」五字并入反白框；正文每四句一列，定宽盒 \hfil 居列中；
% \fontsize{24}{41} 列距对齐全书界行（121+41k）；\clearpage 使述赞独占一页。
\clearpage
\shuzanhead{索隐述赞曰}\linebreak
{\shuzanbody\fontsize{24}{41}\selectfont
\hbox to \hsize{\hfil 清梠凛凛。壮气熊熊。各竭诚义。递为雌雄。\hfil}\linebreak
\hbox to \hsize{\hfil 和璧聘返。渑池好通。负荆知惧。屈节推工。\hfil}\linebreak
\hbox to \hsize{\hfil 安边定策。颇。牧之功。\hfil}%
\par}
\end{document}
